% Kết luận. Mọi con số dưới đây đã xuất hiện ở mục Kiểm tra tính đúng đắn hoặc mục Kết quả và phân tích.

Nhóm cài đặt ba spinlock, một sense-reversing barrier và ba họ concurrent counter, kiểm tra chúng bằng
stress test, ThreadSanitizer và mutation testing, rồi đo chúng trên \PCname{apple-m5} và
\PCname{LAP1184}. Câu trả lời cho RQ1 tới RQ5 dựa trên \PCname{apple-m5}. RQ6 so hai máy.

RQ1. Mọi cài đặt qua mọi bài kiểm tra, và bộ kiểm tra bắt được cả 14 biến thể cố ý làm hỏng. Không kỹ
thuật nào tự nó đủ. Lỗi thuật toán lộ ra dưới stress test, còn lỗi memory order như nhả lock bằng store
\texttt{relaxed} chỉ ThreadSanitizer bắt được.

RQ2. Tới 4 thread, cả ba spinlock bám cận dưới lý tưởng trong khoảng 1,24 lần. Ở 20 thread, TTAS chậm
2,48 lần cận dưới và TAS 3,48 lần, đúng thứ tự mà lưu lượng coherence dự đoán. MCS sụp ngay khi số thread
bằng số core, vì nó trao lock theo thứ tự cố định. Mỗi lần thread kế tiếp không có core là một lượt lập
lịch bị mất. Lock thư viện không sụp vì thread chờ ngủ trong kernel.

RQ3. Backoff giúp TAS nhiều nhất. Ở 20 thread, TAS với sleep backoff chỉ tốn 0,31 lần thời gian của TAS.
Nhưng nó đạt điều đó bằng cách cho các thread chạy lần lượt, nên độ công bằng đúng bằng $1/p$. Lock nhanh
nhất là lock kém công bằng nhất. MCS công bằng tuyệt đối chừng nào mỗi thread còn có core.

RQ4. Chi phí một episode của barrier tăng khoảng 53~ns cho mỗi thread trong cụm P-core, dốc hơn khi
thread tràn sang E-core, và tăng khoảng 2\,000 lần khi số thread vượt số core. Spin thuần là lựa chọn tốt
khi mỗi thread có core riêng. Khi vượt số core, chờ bằng \texttt{sched\_yield} nhanh nhất.

RQ5. Không counter nào mở rộng. Fetch-and-add nhanh hơn vòng CAS 2,1 tới 2,7 lần từ 4 thread vì số lần
CAS hỏng trên mỗi lần tăng tăng từ 0,09 ở 2 thread lên 1,59 ở 8 thread. Backoff đưa số lần hỏng về gần 0
nhưng tuần tự hoá các thread. Mutex thư viện thắng fetch-and-add khi không có việc giữa hai lần tăng,
vì người giữ lock tăng nhiều lần liên tiếp trên cache line nằm sẵn trong cache của mình.

RQ6. Một phần. Trên \PCname{LAP1184} (x86-64, Linux trong máy ảo WSL2), mọi cài đặt cũng qua mọi bài
kiểm tra, và các thứ hạng chính vẫn giữ: TTAS nhanh hơn TAS, MCS sụp khi mọi luồng phần cứng đều bận,
fetch-and-add nhanh hơn vòng CAS. Ba khác biệt đến từ nền tảng. Với TSO, \texttt{seq\_cst} không tốn
thêm gì (0,75 tới 1,14 lần \PCname{cas}), và stress test không bao giờ bắt được lỗi memory order. Barrier
spin sụp ngay ở 12 thread, bằng số luồng phần cứng, có lẽ vì Windows phải lấy core của các CPU ảo.
\texttt{sleep\_for} ngắn nhất kéo dài khoảng 61~\textmu s, gấp 23 lần trên \PCname{apple-m5}, nên các
biến thể ngủ chậm hơn nhiều.

Bảy dự đoán viết trước khi đo bị số liệu bác bỏ toàn phần hoặc một phần. MCS không phải spinlock nhanh
nhất khi số thread chưa vượt số core (D2.1). Lock chỉ bão hoà tại $p = 1 + \mathrm{OW}$ khi
$1 + \mathrm{OW} \le 4$ (D2.3), và sleep backoff vẫn giúp TAS ở nhiều thread dù có việc ngoài
critical section (D3.2). Chi phí barrier không tăng theo một độ dốc duy nhất (D4.1).
Mutex thư viện và vòng CAS ở 2 thread nhanh hơn fetch-and-add (D5.1). \texttt{seq\_cst} làm CAS chậm
khoảng 2 lần (D5.4). Khi vượt số core, MCS mất cả tính công bằng (D3.3). Điểm chung của các dự đoán sai
là chúng chỉ tính lưu lượng coherence mà bỏ qua scheduler của hệ điều hành và việc máy có hai loại core.

Từ các kết quả này, nhóm khuyến nghị ba điều. Khi số thread có thể vượt số core, dùng lock thư viện
hoặc TTAS có backoff, và chỉ dùng MCS khi mỗi thread chắc chắn có core riêng. Dùng barrier spin khi số
thread không vượt số core, và barrier có \texttt{sched\_yield} khi có thể vượt. Dùng fetch-and-add cho
counter, và chia counter theo thread khi nó nằm trên đường nóng. Việc chưa làm gồm phiên bản phân tán bằng MPI-3 RMA
(phần 4 của đề), số đo trên các máy còn lại của nhóm, và một máy x86 chạy Linux trực tiếp để tách ảnh
hưởng của máy ảo.
